% !TEX root = ../nc-cubic-n-fold.tex
\section{点の平行移動と比較核の Picard 類}
\label{sec:nfm-kernel-point-class}

最高次線束型の比較核について、スカラー障害は次数零 Picard 類
$m_P$ の表現障害である。
本節では楕円曲線の rotation の共役式を求め、
可換特殊ファイバーを含む非可換 active 商の相対 THCR/window1 比較を構成する。
これにより、固定入力と良い一ブロックの仮定の下で、
点の平行移動と $m_P$ を結ぶ追加型の非空性とスカラー障害の具体式を証明する。
代数的な active 同値核の assignment や、full 核の同型類の選択を
無条件に構成するものではない。

\subsection{平行移動・線束と rotation の共役}

楕円曲線の点 $b$ を次数零線束 $P_b=\mathcal O_E(b-0)$ と同一視する。
以下の式はまず代数閉体上で証明し、必要な等式を忠実平坦降下して用いる。
右加群の twisted homogeneous coordinate ring の規約は
\cite[\S3, equation~(3.1)]{NFMStaffordVdB}に合わせる。
すなわち $\sigma$ が点加群の shift を表すとき、
coherent sheaf 上の次数 shift は
\begin{equation}
\label{eq:nfm-right-thcr-shift}
 s=\sigma_*\circ(-\otimes L)
\end{equation}
である。特に $s(k(x))=k(\sigma(x))$ となる。
$\sigma^*$ と取り違えると、以下の平行移動の符号が逆になる。

\begin{lemma}[楕円曲線の rotation の共役式]
\label{lem:nfm-elliptic-rotation-conjugacy}
$E$ を標数零の代数閉体上の楕円曲線、$L$ を次数三の線束とする。
\[
 U_{a,b}=t_a^*\circ(-\otimes P_b),\qquad
 \rho=T_{\mathcal O_E}\circ(-\otimes L)
\]
と置く。関手の同型類について
\begin{equation}
\label{eq:nfm-elliptic-rotation-conjugacy}
 \rho U_{a,b}\rho^{-1}\simeq U_{-2a-b,\,3a+b}
\end{equation}
が成り立つ。
$\sigma=t_p$ とし、右加群の規約に従って
$\rho_p=T_{\mathcal O_E}\circ\sigma_*\circ(-\otimes L)$ と置くと、
\begin{equation}
\label{eq:nfm-elliptic-point-intertwiner}
 U_{a,b}\rho_pU_{a,b}^{-1}\simeq\rho
 \quad\Longleftrightarrow\quad
 3a=0,\qquad b=p.
\end{equation}
特に $\rho$ と可換する $U_{a,b}$ は $U_{a,0}$、$a\in E[3]$ に限られ、
これらは次数零 Picard 類を変えない。
\end{lemma}

\begin{proof}
平行移動は次数零 Picard 類に自明に作用するので、
$U_{a,b}$ の同型類の積は $(a,b)$ の加法である。
$t_a^*L\otimes L^{-1}\simeq P_{-3a}$ より
\[
 \operatorname{Ad}_{-\otimes L}(a,b)=(a,b+3a).
\]
また
\[
 \operatorname{Ad}_{T_{\mathcal O_E}}(a,b)=(a-b,b)
\]
である。後式の graph と線束を確認する。
twist の Fourier--Mukai 核は
$\operatorname{Cone}(\mathcal O_{E\times E}\to\mathcal O_\Delta)
=\mathcal O_{E\times E}(-\Delta)[1]$ であり、
$T_{\mathcal O_E}^{-1}k(x)=\mathcal O_E(x)$ となる。
従って共役 $T_{\mathcal O_E}U_{a,b}T_{\mathcal O_E}^{-1}$ は
$k(x)$ を $k(x-a+b)$ へ送り、$\mathcal O_E$ を $P_b$ へ送る。
点層を shift なしの点層へ送る同値の核は、その台である graph 上の線束である。
この場合の graph は $t_{a-b}$ の pullback、線束は $P_b$ なので、所要の式を得る。
これは上記の twist 核を使った畳込み計算とも一致する。

二つの共役作用の積は
\[
 A=\begin{pmatrix}-2&-1\\3&1\end{pmatrix}
\]
であり、\eqref{eq:nfm-elliptic-rotation-conjugacy}を与える。
平行移動は $\mathcal O_E$ を保つので、その pullback は
$T_{\mathcal O_E}$ と可換する。
従って $\rho_p=U_{-p,0}\rho$ であり、intertwining の条件は
\[
 (I-A)\binom ab=\binom p0,
 \qquad\text{すなわち}\qquad
 3a+b=p,\quad -3a=0
\]
に等しい。これで全ての主張が従う。
\end{proof}

この補題は、全自己同値群の分類ではない。
平行移動と次数零線束の tensor からなる部分について、
rotation を保ったまま Picard 類を調整できる範囲を決定している。
特に、この部分の自己同値だけを用いて任意の $m_P$ を零にすることはできない。

\begin{lemma}[右楕円座標環の window1 と内部 shift]
\label{lem:nfm-twisted-elliptic-window}
$A=B(E,L,\sigma)$ を右 twisted homogeneous coordinate ring とする。
$\operatorname{qgr}A\simeq\operatorname{coh}E$ の下で
$\pi A\leftrightarrow\mathcal O_E$、$(1)\leftrightarrow s$ は
\eqref{eq:nfm-right-thcr-shift}に従うものとする。
$A$ に Gorenstein parameter 零の Orlov 同値が適用でき、
window1 を $\Phi_1=q_{\mathrm{sg}}R\Gamma_{\ge1}$ と取る。
このとき
\begin{equation}
\label{eq:nfm-twisted-active-window}
 \kappa_A\longleftrightarrow\mathcal O_E[1],\qquad
 \Phi_1^{-1}(1)\Phi_1=T_{\mathcal O_E}\circ s.
\end{equation}
\end{lemma}

\begin{proof}
この window の存在については
\cite[Theorem~2.5(iii)]{NFMOrlovSingularities}を用いる。
$m>0$ に対して $s^m\mathcal O_E$ は次数 $3m$ の線束なので、
その高次 cohomology は消える。
従って $R\Gamma_{\ge1}\mathcal O_E=A_{\ge1}$ であり、
$0\to A_{\ge1}\to A\to k\to0$ を特異圏へ写すと第一式を得る。
任意の $F$ に対し
\[
 (R\Gamma_{\ge1}F)(1)=R\Gamma_{\ge0}(sF)
\]
であり、次数零の切断は三角形
\[
 R\Gamma_{\ge1}(sF)\longrightarrow R\Gamma_{\ge0}(sF)
 \longrightarrow R\operatorname{Hom}(\mathcal O_E,sF)\otimes k
 \longrightarrow
\]
を与える。
特異圏へ写し、第一式を用いて三角形を回転すると、評価写像の三角形
\[
 R\operatorname{Hom}(\mathcal O_E,sF)\otimes\mathcal O_E
 \longrightarrow sF\longrightarrow\Phi_1^{-1}(1)\Phi_1(F)
 \longrightarrow
\]
を得る。この cone は $T_{\mathcal O_E}(sF)$ である。
従って剰余体と内部 shift の双方について同じ window の規約を得た。
この切断の議論は $A$ の可換性を用いない。
\end{proof}

\subsection{可換特殊ファイバーを含む相対楕円座標環と window1}

以下では $k$ を標数零の体、$R=k[[h]]$、$K=k((h))$ とする。
$g\in k[x,y,z]_3$ は滑らかな平面三次を定め、
\begin{equation}\label{eq:nfm-active-dvr-algebra}
 D_R=R\langle x,y,z\rangle/
 \bigl([y,z]-h\operatorname{sym}g_x,
       [z,x]-h\operatorname{sym}g_y,
       [x,y]-h\operatorname{sym}g_z\bigr)
\end{equation}
の各次数が有限自由で、剰余が $k[x,y,z]$ であると仮定する。
ここで $\operatorname{sym}$ は二次単項式の対称化である。
$\omega\in Z(D_R)_3$ は $\omega\bmod h=g$ を満たすものとし、
$A_R=D_R/(\omega)$ と置く。
この次数別自由性と指定中心元は、支配的な一ブロック族の active 制限で用いる PBW 入力である。

\begin{proposition}[指定中心商の相対 THCR 表示]
\label{prop:nfm-relative-active-thcr}
上の仮定の下で
\[
 J_g=g+\frac{h^2}{24}\det\operatorname{Hess}(g),\qquad
 E_R=V(J_g)\subset\mathbf P^2_R,\qquad L=\mathcal O_{E_R}(1)
\]
と置く。$E_R$ は滑らかな固有 genus-one 曲線族である。
右点加群の関係は、閉ファイバーで恒等となる自己同型
$\sigma:E_R\xrightarrow{\sim}E_R$ の graph を含み、
その一般ファイバーには固定点がない。
さらに
\begin{equation}\label{eq:nfm-relative-active-thcr}
 A_R\xrightarrow{\sim}
 B(E_R,L,\sigma)
 =\bigoplus_{m\ge0}H^0(E_R,L_m),\qquad
 L_m=\bigotimes_{j=0}^{m-1}\sigma^{j*}L
\end{equation}
が次数一の座標切断から定まり、
\[
 \operatorname{qgr}A_R\simeq\operatorname{coh}E_R,
 \qquad (1)\longleftrightarrow s=\sigma_*\circ(-\otimes L)
\]
となる。
代数同型の各次数は $R$-有限自由加群の同型であり、
剰余体・分数体への底変更後にも同じ点対応と THCR 表示を与える。
\end{proposition}

\begin{proof}
$u=(x,y,z)^{\mathsf t}$、$H(u)=\operatorname{Hess}(g)(u)$ とし、
$[u]_\times v=u\times v$ と置く。
右点関係の行列は
\[
 M_h(u)=[u]_\times-\frac h2H(u)
\]
である。対称行列 $H$ と三次の Euler 恒等式から
\begin{equation}\label{eq:nfm-active-point-matrix-determinant}
 \det M_h(u)=-\frac h2u^{\mathsf t}H(u)u
               -\frac{h^3}{8}\det H(u)=-3hJ_g(u).
\end{equation}
$J_g\bmod h=g$ なので、各次数の乗法 $J_g$ は剰余で単射である。
単元小行列式によってその余核は $R$-自由となり、$E_R$ は $R$-flat である。
滑らかでない locus の $\operatorname{Spec}R$ への像は proper 性により閉であり、
閉点を含まないので空である。これで滑らかな genus-one 族を得る。

$E_R$ 上で $\det M_h=0$ であり、閉ファイバーの $M_0(u)=[u]_\times$ は
各点で階数二である。同じ proper 性により階数一以下の locus は空であり、
$\ker M_h$ は直線部分束となる。その直線が定める射を
$\sigma:E_R\to\mathbf P^2_R$ と書く。
三次の偏微分の対称性より
\[
 M_{-h}(v)u=-M_h(u)v.
\]
従って $v=\sigma(u)$ に対し $hJ_g(v)=0$ であり、
$E_R$ の $R$-flatness から $J_g(v)=0$ である。
同じ構成を $-h$ で行うと逆射を得るため、$\sigma$ は $E_R$ の自己同型である。
$\ker M_0(u)=ku$ なので閉ファイバーで恒等となる。
一般ファイバーで $\sigma(u)=u$ なら
$M_h(u)u=-h\nabla g(u)=0$ となり、$V(g)$ の滑らかさに反する。

忠実平坦な有限非分岐拡大後、閉ファイバーの閉点を持ち上げて $E_R$ に原点を選べる。
原点を保つ楕円曲線自己同型の非分岐性から、
$\sigma\bmod h=\mathrm{id}$ は $\sigma=t_p$、$p\bmod h=0$ を意味する。
$E_R[3]$ は有限 \'etale なので、$3p=0$ なら $p=0$ となる。
一般ファイバーに固定点がないことから $p\ne0$ であり、
従って $\sigma_K^3\ne\mathrm{id}$ である。この不等式は原点の選択に依存しない。

$m>0$ について $\deg L_m=3m$ なので、$H^1(E_R,L_m)=0$、
$H^0(E_R,L_m)$ は階数 $3m$ の自由 $R$-加群で、形成は底変更に両立する。
積を $a\cdot b=a\otimes\sigma^{m*}b$（$\deg a=m$）と定める。
次数一の座標切断は上の graph 上の三つの関係を満たすので、
$D_R\to B(E_R,L,\sigma)$ を与える。

指定した $\omega$ が零に写ることを確認する。
一般ファイバーの関数体上で $L$ を有理的に自明化すると、
THCR は $K(E_K)[z;\sigma^*]$ に埋め込まれる。
$\omega$ の像が非零な $fz^3$ なら、次数一の切断 $az,bz$ との可換性から
\[
 f\sigma^{3*}(a)=a\sigma^*(f),\qquad
 f\sigma^{3*}(b)=b\sigma^*(f),\qquad
 \sigma^{3*}(a/b)=a/b
\]
を得る。$L$ は平面埋込みを与えるので、このような比は $K(E_K)$ を生成する。
従って $\sigma_K^3=\mathrm{id}$ となって矛盾する。
THCR の次数三は $R$-torsionfree なので、$\omega$ の像は $R$ 上でも零である。
この議論では中心三次元の一般公式やその scalar 規格化を用いない。

$\omega\bmod h=g$ の乗法は単射である。
再び各次数の単元小行列式により、$\omega$ は正則で、
$(A_R)_m$ は $m=0$ で階数一、$m>0$ で階数 $3m$ の自由 $R$-加群となる。
得られた $A_R\to B(E_R,L,\sigma)$ の剰余は
\[
 k[x,y,z]/(g)\xrightarrow{\sim}
 \bigoplus_{m\ge0}H^0(E_0,\mathcal O_{E_0}(m))
\]
という通常の平面三次の座標環の同型である。
次数ごとの Nakayama と等しい自由階数から
式\eqref{eq:nfm-relative-active-thcr}を得る。

相対的な ample pair の条件も確認する。
任意の coherent 層 $\mathcal F$ に対し $\mathcal G_m=\mathcal F\otimes L_m$ と置く。
曲線族の cohomological dimension は一なので、
\[
 0\to\mathcal G_m[h]\to\mathcal G_m\xrightarrow h h\mathcal G_m\to0,
 \qquad
 0\to h\mathcal G_m\to\mathcal G_m\to\mathcal G_m/h\mathcal G_m\to0
\]
から
\[
 H^1(E_R,\mathcal G_m)/hH^1(E_R,\mathcal G_m)
 \simeq H^1(E_0,(\mathcal F/h\mathcal F)\otimes\mathcal O_{E_0}(m))
\]
となる。右辺は $m\gg0$ で零であり、左辺の $H^1$ は有限 $R$-加群なので零となる。
同じ完全列から $H^0(\mathcal G_m)\to H^0(\mathcal G_m/h\mathcal G_m)$ は全射となる。
閉ファイバーの大きな twist の生成切断を持ち上げ、
評価写像の余核の proper support を用いれば $\mathcal G_m$ 自身も大域生成される。
$s^{-m}\mathcal O=L_m^{-1}$ なので、これらは $(\mathcal O,s)$ の ample-pair 条件である。
$\sigma^{-m*}L_m$ も剰余が $\mathcal O_{E_0}(m)$ であり、同じ議論は反対側にも適用できる。
一般 noetherian 底上の ample-pair 定理
\cite[Proposition~7.2 and Lemma~7.6]{NFMKeelerAmple}により qgr 同値を得る。
右加群での shift は\cite[Section~3, equation~(3.1)]{NFMStaffordVdB}の規約で
$s=\sigma_*\circ(-\otimes L)$ である。
\end{proof}

\begin{lemma}[三変数の相対 Gorenstein 条件]
\label{lem:nfm-active-dvr-gorenstein}
前命題の $D_R,A_R$ は左右 noetherian であり、
\[
 \operatorname{gldim}D_R\le4,\qquad
 R\underline{\operatorname{Hom}}_{D_R}(R,D_R)\simeq R(3)[-3],
 \qquad
 R\underline{\operatorname{Hom}}_{A_R}(R,A_R)\simeq R[-2]
\]
が左右で成立する。$A_R$ の左右 injective dimension は高々三であり、
Burke--Stevenson の意味で parameter 零の Gorenstein 環で、条件 $\chi$ を満たす。
これらの有限相対分解は剰余体と分数体への底変更に両立する。
\end{lemma}

\begin{proof}
入力B\cite{NFNCCubicB}の \texttt{local\_model.tex}、
補題 \texttt{lp:lem:noetherian}, \texttt{lp:lem:resolution} と
命題 \texttt{lp:prop:koszul}, \texttt{lp:prop:gorenstein} の次数別の証明を用いる。
三変数への適用は同入力 \texttt{explicit\_example.tex} の
補題 \texttt{gre:lem:elliptic-mf} にも明記されている。
必要な点を本設定で確認する。

$h$-随伴次数環は中央多項式拡大 $k[x,y,z,t]$ の商なので noetherian である。
斉次イデアルの記号の有限生成元を持ち上げ、
各内部次数の有限 $R$-加群に Krull 交叉定理を用いると $D_R$ の noetherian 性を得る。
任意のイデアルは内部次数の最高次項による有限降下で扱える。
多項式環の augmentation と対角の Koszul 分解を、各次数の分裂全射を用いて持ち上げる。
各段階の核は次数ごとに自由であり、最後の核は Nakayama によって零となる。
得られる二つの長さ三の分解の項は、それぞれ
$D_R(-j)^{\binom3j}$ と $D_R^e(-j)^{\binom3j}$ である。
対角分解と $\operatorname{gldim}R=1$ から $\operatorname{gldim}D_R\le4$ を得る。

augmentation 分解の双対を各内部次数で調べる。
剰余の Koszul 双対複体は内部次数 $-3$、cohomological degree 三以外で完全であり、
有界有限自由複体の単元成分の消去により、この完全性は $R$ 上へ持ち上がる。
内部次数 $-3$ の最後の項には $R$ だけが残るので、最初の導来 Hom 公式を得る。
中心正則列
$0\to D_R(-3)\xrightarrow\omega D_R\to A_R\to0$ による change of rings は
\[
 R\underline{\operatorname{Hom}}_{D_R}(M,D_R)
 \simeq R\underline{\operatorname{Hom}}_{A_R}(M,A_R)(3)[-1]
\]
を与える。$M=R$ から第二式、一般の $M$ から injective dimension の評価を得る。
全て左右で同じであり、$A_R$ は noetherian な $D_R$ の商である。
$R$-flatness と\cite[Definition~5.1 and Lemma~6.2]{NFMBurkeStevenson}から条件 $\chi$ が従う。
分解の各内部次数は $R$ 上で分裂しているので、述べた底変更にも両立する。
\end{proof}

\begin{proposition}[source の相対 window1 と二つの体への底変更]
\label{prop:nfm-relative-active-window-basechange}
命題\ref{prop:nfm-relative-active-thcr}の仮定の下で、
window1 と余核は $R$-線型 dg Morita 同値
\begin{equation}\label{eq:nfm-relative-active-window-proved}
 \operatorname{MF}^{\mathbf Z}(D_R,\omega)
 \simeq D^b_{\mathrm{sg,dg}}(\operatorname{gr}A_R)^{\mathrm{idem}}
 \simeq\operatorname{Perf}(E_R)
\end{equation}
を与える。これらは $R\to k$ および $R\to K$ の導来底変更に両立し、
augmentation stabilization と内部 shift はそれぞれ
\[
 \kappa_{A_R}\longleftrightarrow\mathcal O_{E_R}[1],\qquad
 (1)\longleftrightarrow
 T_{\mathcal O_{E_R}}\circ\sigma_*\circ(-\otimes L)
\]
となる。指定周期は選んだ $\omega$ の余核二拡張から移し、任意の scalar 倍で取り替えない。
\end{proposition}

\begin{proof}
前補題は\cite[Theorem~6.4(3)]{NFMBurkeStevenson}の全ての環論的仮定を満たす。
従って
\[
 \Phi_{1,R}=q_{\mathrm{sg}}R\Gamma_{\ge1}:
 D^b(\operatorname{coh}E_R)\longrightarrow D^b_{\mathrm{sg}}(\operatorname{gr}A_R)
\]
は同値である。
$R\Gamma_{\ge1}$ が有界有限生成の対象に値を持つことは、
同論文 Proposition~4.18 とその直前の有限生成性の議論による。
有限 cohomological dimension と $\chi$ は同論文 Lemmas~6.1--6.2 による。
derived section、内部次数切断、dg 商を合成した標準 enhancement の quasi-functor を用い、
全ての shift に対する完全忠実性から dg Morita 同値を得る。
$E_R$ は正則な二次元 scheme なので $D^b(\operatorname{coh}E_R)=\operatorname{Perf}(E_R)$ である。
余核による最初の dg 比較は入力Bの
\texttt{gre:lem:elliptic-mf} と \texttt{tw:lem:stabilization} の
exact 圏・有限 totalization による構成を使う。
その仮定である ambient の有限大域次元、中心元の正則性、商の有限 injective dimension は全て前補題で確認済みである。

底変更を関手そのものについて確認する。
$F\in\operatorname{Perf}(E_R)$ に対し、derived section の内部次数 $m$ 成分は
\[
 (R\Gamma_*F)_m=R\Gamma(E_R,s^mF)
   \simeq R\Gamma(E_R,F\otimes L_m)\quad(m\ge0)
\]
である。後者は $R$-perfect であり、その自然な比較写像は任意の底変更で擬同型となる
\cite[Lemma~36.30.4, Tag~0B91]{NFMStacksPerfectBasechange}。
この比較は切断の積に自然なので、$\ell=k,K$ について
\begin{equation}\label{eq:nfm-active-window-section-basechange}
 (R\Gamma_{\ge1}F)\otimes_R^{\mathbf L}\ell
 \xrightarrow{\sim}
 R\Gamma_{\ge1}(F\otimes_R^{\mathbf L}\ell)
 \quad\text{in }D(\operatorname{Gr}A_\ell)
\end{equation}
を得る。これは内部次数ごとの擬同型を組み合わせたものであり、
無限個の内部次数を有限個に切り捨てる議論ではない。
$F$ の固定された有限振幅と曲線族の cohomological dimension 一により、
cohomological な振幅は内部次数に一様である。
有限生成性は先に確認した Proposition~4.18 による。

$-\otimes_R^{\mathbf L}\ell$ は Tor 振幅高々一であり、
有界有限生成 graded 複体を同種の複体へ、perfect $A_R$-複体を perfect $A_\ell$-複体へ送る。
従って特異 dg 商への実際の底変更関手があり、
式\eqref{eq:nfm-active-window-section-basechange}は $\Phi_1$ の底変更図式を与える。
有限自由 MF の行列は、閉ファイバーで積が非零多項式 $g$ の乗法となるため単射である。
各内部次数で単元小行列式を使えば、その余核は $R$-自由である。
よって余核、周期分解、有限 totalization も $k,K$ への底変更と可換する。

これで二つの体上の比較との可換図式を、係数変更関手から直接構成した。
曲線側では $\operatorname{Perf}(E_R)\otimes_R^{\mathbf L}\ell
\simeq\operatorname{Perf}(E_\ell)$ が成り立つ。
実際、ample な $L$ の全ての整数冪は $R$ 上に定義され、
それらとその引戻しはそれぞれの perfect 圏を厚生成する。
これらの間の導来 Hom の底変更は、双対との tensor 積へ
Tag~0B91 を適用すれば分かる。
従って自然な係数変更は生成対象間で完全忠実で、ファイバー側も生成する。
各辺の同値と上の図式から、MF 側の自然な係数変更も Morita 同値を与える。
任意の特異圏の dg 商が任意の非平坦底変更に可換する、という一般則は用いていない。

最後に $R\Gamma_{\ge1}\mathcal O_{E_R}=(A_R)_{\ge1}$ と
$0\to(A_R)_{\ge1}\to A_R\to R\to0$ から augmentation の式を得る。
補題\ref{lem:nfm-twisted-elliptic-window}の切断三角形を $R$ 上で用いれば、
内部 shift は相対評価写像の cone
$T_{\mathcal O_{E_R}}\circ s$ へ移る。
指定周期の比較は、入力B \texttt{tw:eq:period-extension} の余核二拡張をそのまま使う。
\end{proof}

\subsection{恒等から変形した相対比較}

\begin{lemma}[形式的な楕円比較が定める Picard 類]
\label{lem:nfm-formal-elliptic-point-line}
$R$ を剰余標数零の完備 DVR とし、$E/R$ を滑らかな固有 genus-one 曲線、
$L$ を相対次数三の実際の線束とする。
$\sigma=t_p$ は閉ファイバーで恒等となる平行移動とする。
同じ $E/R$ 上に perfect な相対 Fourier--Mukai 核で与えられる同値 $F$ があり、
その核の閉ファイバーは $\mathcal O_\Delta$ とする。
さらに
\begin{equation}
\label{eq:nfm-formal-elliptic-intertwining}
 F\bigl(T_{\mathcal O_E}\circ\sigma_*\circ(-\otimes L)\bigr)F^{-1}
 \simeq T_{\mathcal O_E}\circ(-\otimes L)
\end{equation}
を仮定する。
このとき $F(\mathcal O_E)$ は次数零線束 $M$ であり、
\begin{equation}
\label{eq:nfm-formal-point-line-class}
 [M]=p\quad\text{in }\operatorname{Pic}^0_{E/R}(R)
\end{equation}
となる。
\end{lemma}

\begin{proof}
閉ファイバーの閉点を取り、有限非分岐拡大後にその点を $R$ 上の切断へ持ち上げる。
従って忠実平坦降下のため、原点を持つ楕円曲線の場合を証明すればよい。
核を $Q$ と書く。
第一射影 $E\times_R E\to E$ に関して $Q$ は relative perfect であり、
閉ファイバーの各点上の導来制限は長さ一の点層である。
relative Tor amplitude が $[0,0]$ となる locus は開で、閉ファイバー全体を含む。
その補集合の $\operatorname{Spec}R$ への像は proper 性により閉であり、
閉点を含まないので空である。
従って $Q$ は第一射影について平坦な層で、各ファイバーの長さは一である。
同じことが第二射影についても成り立つ。

長さ一の層の族の台は section であり、その上の層は線束である。
両射影についてこの記述を用いると、$Q$ は相対自己同型 $f$ の graph 上の線束となる。
閉ファイバーで $f$ は恒等、線束は自明であり、cohomological shift は零である。
原点を保つ楕円曲線自己同型の非分岐性から、
閉ファイバーで恒等となる $f$ は平行移動である。
従って $F$ の同型類は $U_{a,b}$ と書け、$a,b$ は閉点で零となる。
線束の次数は局所定数なので零である。

一般ファイバーの代数閉包上で
補題\ref{lem:nfm-elliptic-rotation-conjugacy}を適用すると、
\eqref{eq:nfm-formal-elliptic-intertwining}は $3a=0$、$b=p$ を与える。
これらは相対 Jacobian の切断の等式なので $R$ 上へ降りる。
$E[3]$ は有限 \'etale であり、閉点で零となる $a$ は零である。
従って $F(\mathcal O_E)$ の Picard 類は $p$ となる。
この等式は原点の選択に依存しないため、最初の非分岐拡大からも降下する。
\end{proof}

この補題は、perfect な複体が閉ファイバーで線束となることだけから
線束の Picard 類まで一定であると主張するものではない。
その類を決定する追加条件は rotation の intertwining であり、
一般には $M$ は自明線束ではなく $p$ を表す。

\subsection{点の平行移動を指定した full 比較核の型}

定理\ref{thm:nfm-kernel-active-line-bound}の split 表示、
$\kappa_{\mathrm{nc}}$、$q_Y$、window1 を固定する。
その底を $S$ と書く。
系\ref{cor:nfm-semiclassical-full-residual}の半古典規格化により、
命題\ref{prop:nfm-point-normalized-jacobian-cover}の Jacobian 同型の切断
$\theta_{\mathrm{sc}}:J_E\xrightarrow{\sim}J_C$ が選ばれる。
命題\ref{prop:nfm-point-residual-isom-cover}の平面同型被覆をこの切断で引き戻し、
\[
 \mathfrak I_{\mathrm{sc}}
   =\mathfrak I\times_{\mathfrak D_J,\theta_{\mathrm{sc}}}S
   \longrightarrow S
\]
と置く。これは階数九の有限 \'etale な $J_E[3]$-torsor である。
この被覆上では
\[
 \phi:E_{\mathrm{pt}}\xrightarrow{\sim}C_{\mathrm{res}},\qquad
 \operatorname{Jac}(\phi)=\theta_{\mathrm{sc}},\qquad
 p_{\mathrm{res}}=\theta_{\mathrm{sc}}(p_{\mathrm{pt}})
\]
となる。以下では $\theta=\theta_{\mathrm{sc}}$ と固定する。
$p_{\mathrm{pt}}$ は系\ref{cor:nfm-point-shift-residual-transfer}の切断である。
最高次線束型の核 $P$ に、さらに
\begin{equation}
\label{eq:nfm-kernel-point-line-type}
 m_P=p_{\mathrm{res}}
\end{equation}
を課す。
この条件は $\phi$ の Jacobian 上の作用だけに依存し、
$p_{\mathrm{res}}$ 自体は $S$ 上で定義される。
従って平面同型の三捩れの選択を忘れても条件は定義できるが、
そのことから元の底上の full 核の同型類が選べるとは結論しない。

支配的な完備族に、前小節の相対構成を適用する。
その係数体を $k$、$R=k[[h]]$ とする。
中心一次変数を $U=0$ とした二次関係は、$g=F|_{U=0}$ だけに依存し、
式\eqref{eq:nfm-active-dvr-algebra}に一致する。
中心変数を含まない PBW 単項式が商の基底となるので、
$D_R=B_h/(U)$ は各次数有限自由で剰余は $k[x,y,z]$ である。
選んだ中心三次の像 $\omega=\Omega_h\bmod(U)$ は中心であり、$\omega\bmod h=g$ である。
従って相対点曲線 $E_R=V(J_g)$ と $L=\mathcal O_{E_R}(1)$ に対し、
命題\ref{prop:nfm-relative-active-thcr}と
命題\ref{prop:nfm-relative-active-window-basechange}が適用でき、
\begin{equation}
\label{eq:nfm-relative-elliptic-window-comparison}
 \operatorname{MF}^{\mathbf Z}(B_h/(U),\Omega_h\bmod U)
 \simeq\operatorname{Perf}(E_R)
\end{equation}
を閉ファイバー・分数体への導来底変更と両立して得る。
この比較は augmentation stabilization を $\mathcal O_{E_R}[1]$ へ送り、
内部 shift を $T_{\mathcal O_{E_R}}\circ\sigma_*\circ(-\otimes L)$ へ送る。
閉ファイバーでは可換側と同じ window1 を用いる。
ここで相対 THCR/window1 は追加仮定ではなく、
PBW 自由性・中心元・滑らかな $g$ から証明した帰結である。

\begin{proposition}[点の平行移動を指定した核の型の非空性]
\label{prop:nfm-kernel-point-line-nonempty}
固定入力Bと良い一ブロックの仮定を置く。
このとき、非空開の有限 \'etale 全射上に、平面同型 $\phi$ と
rotation・指定周期を保つ full 比較核 $P$ を取り、
\[
 \operatorname{Jac}(\phi)=\theta_{\mathrm{sc}}
\]
および
最高次線束条件\eqref{eq:nfm-kernel-active-line-type}と
\eqref{eq:nfm-kernel-point-line-type}を同時に満たすようにできる。
従来の三つの Euler 条件\eqref{eq:nfm-kernel-numerical-type}も同時に課せる。
これは少なくとも一つの支配的な被覆上での非空性であり、
全ての平面同型の選択や全ての最高次線束型の核に対する主張ではない。
\end{proposition}

\begin{proof}
命題\ref{prop:nfm-completed-central-restriction}の相対 CF 比較を
active 商へ制限する。
標的の方程式は $b_hJ_g$ であり、underlying curve は $E_R=V(J_g)$ である。
potential の単元倍を扱うときは、既存の比較と同じ指定周期を保つ。
構成\eqref{eq:nfm-relative-elliptic-window-comparison}と可換側の window1 により、
この比較は同じ $E_R$ 上の $R$-線型同値 $F$ となる。
滑らかな固有曲線の perfect 圏の間のこの比較は perfect な相対核で実現される。
選んだ CF 比較は閉ファイバーで恒等であり、二つの window もそこで一致するので、
その核の閉ファイバーは $\mathcal O_\Delta$ である。
内部 shift の適合性と
補題\ref{lem:nfm-twisted-elliptic-window}から
\eqref{eq:nfm-formal-elliptic-intertwining}が成立する。
従って補題\ref{lem:nfm-formal-elliptic-point-line}により
\[
 F(\mathcal O_{E_R})=M_R,\qquad [M_R]=p_{\mathrm{pt}}.
\]

ここから分数体 $\widehat K=k((h))$ へ移る。
命題\ref{prop:nfm-mf-linear-restriction}の graded purity と
全 Taylor 成分における中心係数商との可換性から
\[
 q_Y\Phi_P(\kappa_{\mathrm{nc}})
 \simeq
 \bigoplus_{j=0}^{n-1}
 \rho^{-j}(M_{\widehat K})[j+1]^{\oplus\binom{n-1}{j}}
\]
を得る。
補題\ref{lem:nfm-active-line-rotation}により、その最高次 cohomology は
$M_{\widehat K}$ である。
形式模型と代数的残余三次との一次空間・方程式の同定を選ぶと、
中心作用の保存により active 平面も同定され、平面同型 $\phi$ を得る。
曲線上への pullback と pushforward は Picard 類と平行移動の対応を保つので、
代数的残余曲線の frame では
$m_P=\operatorname{Jac}(\phi)(p_{\mathrm{pt}})$ となる。

この $\phi$ は $\mathfrak I_{\mathrm{sc}}$ に属することも確認する。
active frame では $\phi$ は二つの $J_g$ の恒等同型である。
命題\ref{prop:nfm-point-scale-formal-comparison}の正確な尺度公式により、
その Jacobian 同型から定まる尺度は
\[
 \tau_{\mathrm{pt}}\circ\operatorname{Jac}(\phi)^*
       =\frac{h}{b_h}\,W_h\otimes P_0,
       \qquad W_h|_{U=0}=b_hJ_g
\]
である。支配的な中心族の記法では、その先頭項は
$h(F+G)\otimes P_0$ となる。
命題\ref{prop:nfm-point-normalized-jacobian-cover}は
Jacobian 同型と尺度被覆の切断を一対一に対応させる。
系\ref{cor:nfm-semiclassical-full-residual}の二つの切断のうち、
この先頭項を持つものは $\theta_{\mathrm{sc}}$ に対応するものだけである。
従って形式体上で
$\operatorname{Jac}(\phi)=\theta_{\mathrm{sc}}$ であり、
$\phi$ は $\mathfrak I_{\mathrm{sc}}$ の点を与える。

あとは命題\ref{prop:nfm-kernel-active-line-nonempty}の有限表示の議論を用いる。
元の関数体を $K$ とし、full 核・逆核・単位・余単位・rotation・周期と
平面同型 $\phi\in\mathfrak I_{\mathrm{sc}}$ を
$K$ 上有限生成な整域 $D$ へ降ろす。
ここでは単に $\mathfrak I$ への点としてではなく、
有限表示の $\mathfrak I_{\mathrm{sc}}$ への点として降ろすので、
Jacobian 同型と半古典符号の一致も保存される。
必要な開への縮小により、最高次 cohomology の可逆層性・次数零・底変更適合性を得る。
相対 Picard 類 $m_P$ と $p_{\mathrm{res}}$ は相対 Jacobian の切断である。
形式体への忠実平坦拡大後に両者が等しいので、
$\operatorname{Frac}(D)$ 上で既に等しい。
相対 Jacobian の分離性により、この等式は両切断が定義された非空開上に成立する。
従ってその開の閉点へ特殊化すれば、有限分離拡大 $K'/K$ 上で
所要の追加型を満たす full 核と $\phi$ を得る。
元の表示 chart の非空開へ広げ、正規化と分岐 locus の除去を行えば、
有限 \'etale 被覆上の主張となる。
三つの perfect mixed Hom を同じ有限データに含めれば、
従来の Euler 条件も保存される。

この特殊化で保持したのは full 核とその最高次線束の Picard 等式である。
形式的な active 比較核そのものや、その coherent な中心線型性を
代数的に降下させたという主張は用いていない。
\end{proof}

\begin{corollary}[選択済みの核の障害を測る点の平行移動]
\label{cor:nfm-kernel-point-picard-obstruction}
split 表示と上記の $\theta=\theta_{\mathrm{sc}}$ を固定し、
\eqref{eq:nfm-kernel-point-line-type}を満たす full 核の同型類が
既に整合的に選ばれているとする。
$C=C_{\mathrm{res}}$ とし、その実現 gerbe の類を $\alpha$、
次数零 Picard 類の表現障害を
$\delta_C(m)=m^*\beta_{\mathrm{Pic},C}$ と書くと、
\begin{equation}
\label{eq:nfm-kernel-point-brauer-class}
 \alpha=\delta_C(p_{\mathrm{res}}).
\end{equation}
従って $p_{\mathrm{res}}$ が実際の次数零線束で表されれば、比較核は降下する。
また定理\ref{thm:nfm-kernel-tripling-cover}の三分点被覆を
$p_{\mathrm{res}}$ に適用すれば、その階数九の有限 \'etale 被覆上で
最高次線束を規格化した full 比較核を得る。
$p_{\mathrm{res}}\in3\operatorname{Jac}(C/S)(S)$ なら
系\ref{cor:nfm-kernel-picard-mod-three}により $\alpha=0$ である。
\end{corollary}

\begin{proof}
系\ref{cor:nfm-kernel-picard-obstruction}により
$\alpha=\delta_C(m_P)$ であるから、
\eqref{eq:nfm-kernel-point-line-type}を代入すればよい。
三分点被覆と規格化については引用した定理と系を用いる。
\end{proof}

\begin{remark}[適用範囲と残る問題]
最高次線束型の全ての核について $m_P=p_{\mathrm{res}}$ を示したわけではない。
また、追加型の非空性と、元の底でその full 核の同型類を整合的に選ぶことは別である。
本節では source の相対 THCR/window1 と必要な二つの体への底変更を構成したが、
表示変更に対する型の内在性、full 核の同型類の整合、
および代数的 active 核の構成はなお別の問題である。
一方、$\theta$ の符号は半古典規格化により既に選ばれており、
本節の非空な型はその $\theta_{\mathrm{sc}}$ に一致する。
既に選ばれた追加型に対する\eqref{eq:nfm-kernel-point-brauer-class}は
Picard 障害を具体的な点へ評価した式であり、代数的 active 核の存在を前提としない。

固定点のない $\sigma$ については $p_{\mathrm{pt}}\ne0$ である。
従ってこの追加型の最高次線束は幾何的に非自明であり、
$m_P=0$ という十分条件だけでその核を降下させる方針は適切でない。
実際の線束による $p_{\mathrm{res}}$ の表現、または三分可能性を調べる必要がある。
ここで点の平行移動を full Ku の rotation 自体と同一視しているのではなく、
構成した active window における rotation 共役を利用している。
\end{remark}
